\chapter*{TÓM TẮT}
\addcontentsline{toc}{chapter}{TÓM TẮT}

Báo cáo này xây dựng \textbf{ViVLM-nano} --- một mô hình ngôn ngữ--thị giác
(vision--language model, VLM) tiếng Việt cỡ nhỏ mà \emph{toàn bộ phần ngôn ngữ
được tiền huấn luyện từ token đầu tiên}, thay vì mượn một mô hình ngôn ngữ huấn
luyện sẵn như các VLM tiếng Việt hiện có (LaVy, Vintern-1B). Hệ thống đi trọn
vòng đời huấn luyện của một mô hình sinh hiện đại trên đúng một GPU 24\,GB:
(i)~\emph{tiền huấn luyện} một mô hình decoder-only 101 triệu tham số kiến trúc
Llama-style (RoPE, RMSNorm, SwiGLU) trên 3,02 tỷ token tiếng Việt lọc chất lượng
từ FineWeb2-HQ, với bộ tách từ BPE 20.480 mục tự huấn luyện;
(ii)~\emph{hợp nhất thị giác} kiểu LLaVA: bộ mã hóa SigLIP đóng băng, 196 patch
được nén còn 49 token ảnh bằng pixel-shuffle rồi chiếu vào không gian embedding
của mô hình ngôn ngữ, tinh chỉnh có giám sát (SFT) hai pha trên khoảng 77 nghìn
mẫu mô tả ảnh và hỏi--đáp tiếng Việt (UIT-ViIC, KTVIC, OpenViVQA) với loss chỉ
tính trên phần trả lời; (iii)~\emph{học tăng cường} SCST tối ưu trực tiếp CIDEr
với baseline là chính mô hình giải mã tham lam.

Về đo lường, chất lượng ngôn ngữ được theo dõi bằng perplexity và
bits-per-character (thước đo trung lập với bộ tách từ, cho phép so sánh công bằng
với baseline GPT-2 tiếng Việt huấn luyện sẵn); chất lượng đa phương thức đo bằng
BLEU-4 và CIDEr trên các tập kiểm tra UIT-ViIC, KTVIC và OpenViVQA; hiệu quả của
bộ tách từ đo bằng fertility so với tokenizer của PhoGPT và GPT-2. Toàn bộ pipeline
được phát triển kiểu test-driven với 52 unit test chạy ngoại tuyến, checkpoint
khôi phục chính xác từng bit, và nhật ký GPU ghi mỗi 30 giây làm bằng chứng sử
dụng tài nguyên: giai đoạn tiền huấn luyện duy trì mức sử dụng GPU xấp xỉ 100\%
và 22,2/24,5\,GB VRAM (91\%) với thông lượng ổn định 99--100 nghìn token/giây
nhờ \texttt{bfloat16}, FlashAttention và \texttt{torch.compile}
(so với 34 nghìn token/giây và tràn bộ nhớ ở chế độ eager).
Kết quả định lượng đầy đủ (perplexity, fertility, CIDEr trước/sau SCST) được
trình bày trong chương 4. Cùng với đồ án giữa kỳ (Medical VQA --- nhánh phân
loại) và đồ án cuối kỳ sinh mô tả ảnh (nhánh sinh có điều kiện), báo cáo này
lấp phần còn lại của sơ đồ đa phương thức tổng quát của môn học: hợp nhất hai
luồng token thị giác--ngôn ngữ và chuỗi huấn luyện pretrain--SFT--RL.

\textbf{Từ khóa:} mô hình ngôn ngữ, tiền huấn luyện, VLM, tiếng Việt, SFT, SCST, BPE.
\clearpage
